\section{Evaluation}
\label{sec:evaluation}

This section applies the benchmark to a set of detectors. The goal is
not to crown a winner but to demonstrate what the benchmark measures:
detection accuracy on sink-validated labels, and the robustness drop
that the four adversarial mutators expose.

\subsection{Research questions}
\label{sec:eval:rq}

\begin{description}
  \item[RQ1] How accurately do free static analysis tools classify the
    seven sink-shaped CWE classes on the clean test set?
  \item[RQ2] How much accuracy does each tool lose when the test
    samples are rewritten by a semantics-preserving mutator, and which
    mutator costs the most?
  \item[RQ3] Does the diff-based label filter (Section~\ref{sec:dataset})
    change the picture a detector evaluation paints --- that is, are the
    tools being scored against labels they could plausibly satisfy?
\end{description}

\subsection{Setup}
\label{sec:eval:setup}

\paragraph{Test set.} The held-out test split is 132 samples: 67
vulnerable positives spanning the seven CWE classes and 65 \texttt{safe}
hard negatives. The four single mutators and the composed mutator are
applied to the 67 positives, producing five adversarial variant sets of
67 samples each. The hard negatives are not mutated; they appear only in
the \emph{clean} run, where they measure false positives.

\paragraph{Metric.} Every detector prediction is scored one-vs-rest per
CWE class: a sample labelled CWE-X that the tool flags as CWE-X is a
true positive for that class, a sample of any other label flagged as
CWE-X is a false positive, and so on (Section~\ref{sec:methodology}).
From the per-class counts we compute precision, recall, and F1, and
report the macro-average of F1 across the seven classes. Macro-averaging
gives each class equal weight regardless of support, which matters here
because class sizes in the test set range from 3 (CWE-22) to 32
(CWE-89).

\paragraph{Robustness drop.} For a tool and a mutator, the robustness
drop is the macro-F1 on the clean positives minus the macro-F1 on that
mutator's output. Because the hard negatives are absent from the
mutator variants, the clean baseline for this comparison is recomputed
over the 67 positives only, so the subtraction is like-for-like. The
aggregate robustness number is the mean drop across the four single
mutators.

\paragraph{Detectors.} We report two free, locally-installable static
analysis tools: Bandit 1.9.4 and Semgrep 1.163.0. Semgrep is run with
the \texttt{p/python}, \texttt{p/security-audit}, and
\texttt{p/owasp-top-ten} registry rule packs. Each tool emits findings
that we map to the seven CWE classes; for Bandit the mapping is by rule
identifier rather than by Bandit's own \texttt{issue\_cwe} field, which
mislabels \texttt{eval} as CWE-78 and \texttt{yaml.load} as CWE-20.
LLM detectors are evaluated through the same harness with a fixed
classification prompt; those results are reported in
Section~\ref{sec:eval:llm}.

\subsection{Detection on the clean test set (RQ1)}
\label{sec:eval:clean}

Table~\ref{tab:clean-percwe} gives the per-class F1 of the two static
analysis tools on the clean test set.

\begin{table}[h]
\centering
\caption{Per-CWE F1 on the clean test set. $n$ is the number of
positive test samples in the class.}
\label{tab:clean-percwe}
\begin{tabular}{l r r r}
\hline
CWE & $n$ & Bandit & Semgrep \\
\hline
CWE-22  Path Traversal        &  3 & 0.00 & 0.00 \\
CWE-78  OS Command Injection  &  4 & 0.38 & 0.29 \\
CWE-79  Cross-site Scripting  &  7 & 0.60 & 0.33 \\
CWE-89  SQL Injection         & 32 & 0.80 & 0.06 \\
CWE-94  Code Injection        &  6 & 0.86 & 0.80 \\
CWE-502 Insecure Deserial.    & 12 & 0.80 & 0.73 \\
CWE-918 SSRF                  &  3 & 0.00 & 0.00 \\
\hline
Macro-F1 (positives + safe)   &    & 0.49 & 0.32 \\
\hline
\end{tabular}
\end{table}

Two classes score zero for both tools. Neither Bandit nor Semgrep ships
a path traversal (CWE-22) rule that fires on the test samples, and
neither has a server-side request forgery (CWE-918) rule with usable
CWE metadata; Bandit's closest signal, the \texttt{urllib.urlopen}
check, produced five false positives and no true positives. These zeros
are a real finding, not a harness artifact: the benchmark includes CWE
classes that free static analysis does not cover, and the macro-average
makes that gap visible rather than letting a well-covered class mask it.

The two tools diverge most on SQL injection. Bandit reaches F1 0.80 on
CWE-89; Semgrep reaches 0.06, flagging one of 32 positives. The cause is
the composition of the CWE-89 sub-corpus: most of these samples are
fragments drawn from the \texttt{vudenc} source and lack the
framework-recognisable source-to-sink shape that Semgrep's taint rules
require. Bandit's \texttt{B608} rule, which matches inline SQL string
construction lexically, fires far more readily on such fragments. This
divergence is itself a use of the benchmark: it separates lexical
matching from taint-style reasoning on exactly the inputs where the two
approaches disagree.

The headline macro-F1 of 0.49 (Bandit) and 0.32 (Semgrep) is computed
over all 132 clean samples, including the 65 hard negatives. Restricted
to the 67 positives the same tools score 0.59 and 0.38; the gap is the
F1 cost of false positives raised on safe code, and chiefly reflects
Bandit's low-severity \texttt{subprocess} warnings (\texttt{B603},
\texttt{B607}) firing on benign command-execution code.

\subsection{Adversarial robustness (RQ2)}
\label{sec:eval:robustness}

Table~\ref{tab:robustness} reports macro-F1 on the clean positives and
on each mutator variant, with the robustness drop in the final column.

\begin{table}[h]
\centering
\caption{Macro-F1 by variant, computed over the 67 positives.
DC: dead code, SS: string split, VR: variable rename, WE: wrapper
extraction, Comp: composed. Drop is the mean clean$-$mutator difference
over the four single mutators.}
\label{tab:robustness}
\begin{tabular}{l r r r r r r r}
\hline
Tool & Clean & DC & SS & VR & WE & Comp & Drop \\
\hline
Bandit  & 0.59 & 0.59 & 0.59 & 0.59 & 0.59 & 0.59 & 0.00 \\
Semgrep & 0.38 & 0.38 & 0.38 & 0.38 & 0.38 & 0.38 & 0.00 \\
\hline
\end{tabular}
\end{table}

Both static analysis tools are essentially invariant to the four
mutators: the mean robustness drop is 0.001 for Bandit and 0.004 for
Semgrep, and no single mutator moves macro-F1 by more than 0.01. This is
the expected result and it validates the mutator design from the
negative side. Dead code injection, identifier rename, and wrapper
extraction do not alter the lexical sink tokens that Bandit and Semgrep
match on, so a rule-based detector that ignores everything except those
tokens is unaffected. String splitting does break a contiguous SQL
keyword, but Bandit's \texttt{B608} matches on the call shape
(\texttt{cursor.execute} with a non-literal argument) as well as on
keyword substrings, so it survives. A detector with zero robustness drop
is not robust in any deep sense; it is simply not reading the parts of
the program the mutators rewrite. The mutators are built to discriminate
\emph{learned} detectors, where a non-trivial drop is expected; the
static analysis rows establish the floor.

\subsection{LLM detectors}
\label{sec:eval:llm}

The harness evaluates LLM detectors through the same interface: a fixed
system prompt lists the seven CWE classes, the file contents form the
user turn, and the first line of the response is parsed as the
prediction. A full pass over the 132 clean samples and the five mutator
variants is 467 model calls. These results are pending and will be added
once the LLM evaluation is run; the SAST rows above are the established
baseline against which the LLM robustness drop will be read.

\subsection{Effect of the label filter (RQ3)}
\label{sec:eval:labels}

The numbers above are computed against the diff-filtered labels of
Section~\ref{sec:dataset}, whose audited false-positive rate is 26\%.
Run against the unfiltered labels, every tool's apparent precision falls
by construction: a tool that correctly declines to flag a mislabelled
sample is charged a false negative. The filter therefore does not flatter
the tools --- it removes samples on which a correct detector and the
label disagree because the label is wrong. Reporting detector accuracy
against sink-validated labels is what separates this benchmark from an
evaluation run directly on raw CVE-derived tags, where 30--50\% of the
ground truth is noise and a tool is partly being scored on its agreement
with that noise.
